\documentclass[11pt]{article}

\usepackage[a4paper,margin=1in]{geometry}
\usepackage{cmap}
\usepackage[T1]{fontenc}
\usepackage[utf8]{inputenc}
\input{glyphtounicode}
\pdfgentounicode=1
\usepackage{lmodern}
\usepackage{amsmath,amssymb,amsthm}
\usepackage{booktabs}
\usepackage{enumitem}
\usepackage[unicode,pdfencoding=auto,psdextra]{hyperref}

\hypersetup{
  pdftitle={No Positive Margin from Pencils: Xiao's Lemma, a Rank-r Margin Identity, and the Nef Barrier in Parshin-Moret-Bailly Height Arguments},
  pdfauthor={Lukas Geiger},
  pdfsubject={Pencil arguments for upper bounds on arithmetic self-intersections},
  pdfkeywords={Parshin-Moret-Bailly, arithmetic surfaces, Xiao's lemma, Harder-Narasimhan filtration, Bogomolov inequality, Kodaira fibration},
  bookmarksnumbered=true,
  bookmarksopen=true,
  pdflang={en-US},
  colorlinks=true,
  linkcolor=black,
  urlcolor=blue,
  citecolor=black
}

\clubpenalty=10000
\widowpenalty=10000
\displaywidowpenalty=10000
\emergencystretch=1.5em

\newtheorem{theorem}{Theorem}[section]
\newtheorem{proposition}[theorem]{Proposition}
\newtheorem{lemma}[theorem]{Lemma}
\newtheorem{corollary}[theorem]{Corollary}
\newtheorem*{theoremA}{Theorem A}
\newtheorem*{theoremB}{Theorem B}
\newtheorem*{theoremC}{Theorem C}
\theoremstyle{definition}
\newtheorem{definition}[theorem]{Definition}
\newtheorem{problem}[theorem]{Open problem}
\theoremstyle{remark}
\newtheorem{remark}[theorem]{Remark}

\DeclareMathOperator{\rk}{rk}
\DeclareMathOperator{\Pic}{Pic}
\DeclareMathOperator{\ch}{ch}
\newcommand{\OO}{\mathcal{O}}
\newcommand{\cI}{\mathcal{I}}
\newcommand{\cL}{\mathcal{L}}
\newcommand{\cF}{\mathcal{F}}
\newcommand{\ob}{\overline{\omega}}
\newcommand{\Lb}{\overline{L}}
\newcommand{\degh}{\widehat{\deg}}
\newcommand{\bpf}{\textup{(BPF)}}

\title{No Positive Margin from Pencils:\\ Xiao's Lemma, a Rank-$r$ Margin Identity,\\ and the Nef Barrier in Parshin--Moret-Bailly Height Arguments}
\author{Lukas Geiger\\
Independent researcher, Bernau im Schwarzwald, Germany\\
ORCID: \href{https://orcid.org/0009-0005-7296-1534}{0009-0005-7296-1534}}
\date{Draft v0.1, 7 November 2026}

\begin{document}

\maketitle

\begin{abstract}
Parshin--Moret-Bailly height arguments seek an upper bound for the arithmetic
self-intersection $\ob_P^{\,2}$ of the relative dualizing sheaf on the fibers
of a fibration $\phi\colon X\to Y$ over a curve. A recent implementation of this
strategy (Soul\'e, arXiv:2611.05217, version~1) obtains its key estimate from a
pencil: two sections of an auxiliary line bundle $L$ without common zero,
Minkowski's theorem on successive minima, and a transfer from classical to
arithmetic degrees. The leading coefficient of the resulting inequality is the
geometric margin $M=(-L^2+m\deg\phi_*L)/\omega^2$, where $m$ is the fiber degree
of $L$, and an upper bound for $\ob_P^{\,2}$ requires $M>0$. We show that this
never happens. In rank two, the Koszul resolution of the evaluation map gives the
exact identity $-L^2+m\deg\phi_*L=-\ell(Z)-2m\deg V\le 0$ (Xiao's lemma for
pencils), where $Z$ and $V$ are the finite and vertical parts of the base
locus. In rank $r$, we prove a margin identity in terms of the Bogomolov
discriminant of the syzygy kernel together with a nef lemma in the spirit of
Xiao, which gives $M_r\le 0$ whenever $\phi_*L$ is semistable and strict or
partial negativity in the Harder--Narasimhan-unstable case; the arithmetic side
of the unstable case for $r\ge 3$ is left open. As an application, we show that
Proposition~9\,vi) of arXiv:2611.05217v1 is false: in the Kodaira fibration used
there one has $k=4$, at most two of the four square roots over each exceptional
point are effective, and $\deg\phi_*L=-7k+4k_{\mathrm{eff}}\le -20<-3k$. The step
that fails uses the converse of the implication proved in Proposition~5\,ii) of
that preprint. We add brief remarks on cross terms, on single sections and jets
along curves, and on source bundles not pulled back from the base. No statement
about the $abc$ conjecture is made.
\end{abstract}

\section{Introduction}\label{sec:intro}

Parshin~\cite{Pa88} proposed an arithmetic analogue of the
Bogomolov--Miyaoka--Yau inequality for complex surfaces (for the inequalities of
Bogomolov and Miyaoka see~\cite{Bo78,Mi77}): an upper bound for the arithmetic
self-intersection of the relative dualizing sheaf of an arithmetic surface in
terms of the discriminant of the base. Parshin and
Moret-Bailly~\cite{Pa88,MB90} showed that such a bound, applied to the fibers of
a suitable fibration, has strong Diophantine consequences, including effective
Mordell-type statements and forms of the $abc$ conjecture
(see also~\cite{So05}). We call arguments
of this kind \emph{Parshin--Moret-Bailly (PMB) height arguments}.

In a recent preprint~\cite{So26} (arXiv:2611.05217, version~1, submitted on
3~November 2026), Soul\'e implements this strategy for a Kodaira fibration
$\phi\colon X\to Y$ with fibers of genus~$6$ over an \'etale cover of a curve of
genus~$2$~\cite{Ko67,MD85}. The upper bound for $\ob_P^{\,2}$ on the fiber $X_P$
comes from a \emph{pencil}: an auxiliary line bundle $L$ on $X$, two sections of
$L_P$ without common zero~\cite[Prop.~3]{So26}, Minkowski's theorem on successive
minima, and a transfer from classical to arithmetic degrees~\cite[Cor.~1]{So26}.
The leading coefficient of the resulting inequality is a purely geometric number,
which we call the \emph{margin},
\begin{equation}\label{eq:margin-intro}
M(L)=\frac{-L^2+m\deg\phi_*L}{\omega^2},
\end{equation}
where $m$ is the degree of $L$ on a fiber and $\omega=\omega_{X/Y}$. The argument
bounds $\ob_P^{\,2}$ from above if and only if $M(L)>0$ (Section~\ref{sec:setting}).
In~\cite{So26} the positivity of $M(L)$ rests on the lower bound
$\deg\phi_*L\ge -3k$ of~\cite[Prop.~9\,vi)]{So26}.\footnote{We cite~\cite{So26}
by the numbers printed in version~1. In the \TeX{} source of version~1, the
statements printed as Propositions~3, 4, 5, 8, and~9 carry the labels
\texttt{prop12}, \texttt{prop3}, \texttt{prop13}, \texttt{prop5}, and
\texttt{prop6}, while Lemma~3 carries the label \texttt{lm15}. Equation numbers
such as (24) and (28)--(36) are set by hand in the source and agree with the
printed ones.}

This note shows that the margin of a pencil is never positive and draws the
resulting consequence for~\cite{So26}.

\begin{theoremA}[Theorem~\ref{thm:rank2}]
Let $\phi\colon X\to Y$ be a fibration of a smooth projective surface over a
smooth projective curve with integral fibers, and let $L$ be a line bundle with
$h^0(L_\eta)=2$ whose linear system on the generic fiber is base-point-free. Let
$L(-V)\otimes\cI_Z$ be the image of $\phi^*\phi_*L\to L$. Then
\[
-L^2+m\deg\phi_*L=-\ell(Z)-2m\deg V\le 0 .
\]
Hence $M(L)\le 0$, with equality only if the evaluation map is surjective.
\end{theoremA}

This is Xiao's lemma for pencils~\cite{Xi87} in exact form; no stability
hypothesis is needed.

\begin{theoremB}[Theorem~\ref{thm:rankr}, Corollary~\ref{cor:ss}, Proposition~\ref{prop:unstable}]
For $r=\rk\phi_*L\ge 2$ the relevant margin is
$M_r(L)=(-L^2+\tfrac{2m}{r}\deg\phi_*L)/\omega^2$, and
\[
M_r(L)=-\frac{\Delta(K)+2rm\deg V+2(r-1)\ell(Z)}{r\,\omega^2},
\]
where $\Delta(K)$ is the Bogomolov discriminant of the syzygy kernel $K$. If
$\phi_*L$ is semistable, then $M_r(L)\le 0$. If $\phi_*L$ is unstable, then
$M_r(L)<0$ when the maximal destabilizing subsheaf has no horizontal base divisor,
and in general every coefficient built from two Harder--Narasimhan slopes as in
Proposition~\ref{prop:unstable}\,(b) is nonpositive.
\end{theoremB}

The semistable case rests on a nef lemma in the spirit of Xiao's method: the
$\mathbb{Q}$-divisors $L-B_j-V_j-\mu_jF$ attached to the Harder--Narasimhan
filtration of $\phi_*L$ are nef (Lemma~\ref{lem:nef}). The arithmetic counterpart
of the unstable case for $r\ge 3$ is left open (Open problem~\ref{prob:anomaly}).

\begin{theoremC}[Theorems~\ref{thm:soule-count} and~\ref{thm:soule-structural}]
In the Kodaira fibration of~\cite{So26} one has $k=4$, at most two of the four
square roots of the hyperelliptic class over each exceptional point are
effective, and
\[
\deg\phi_*L=-7k+4k_{\mathrm{eff}}\le -20<-12=-3k ,
\]
where $4k_{\mathrm{eff}}\le 8$ is the number of points of $Y$ over which the
relevant square root is effective. Hence~\cite[Prop.~9\,vi)]{So26} is false, and
the leading coefficient of~\cite[(36)]{So26} is at most $-5/6$. A second route
gives the same conclusion: the generic linear system is base-point-free for every
choice of the component $Y$ and of the square root $\cL$, and Theorem~A yields
$\deg\phi_*L\le -10k/3<-3k$; base-point-freeness is the property used
in~\cite[(28)]{So26} to apply~\cite[Prop.~3]{So26}.
\end{theoremC}

The step that fails in the proof of~\cite[Prop.~9]{So26} is inequality~(24) of
that paper. It uses the implication ``$P_1$ exceptional $\Rightarrow
h^0(\lambda)=1$'', whereas~\cite[Prop.~5\,ii)]{So26} proves the converse
(Remark~\ref{rem:gap}). Section~\ref{sec:scope} records briefly why the same sign
appears for cross terms, for single sections and jets evaluated along curves, and
why a positive term can only come from a bundle that is not pulled back from the
base.

\paragraph{Claim Boundary.}
This note makes no statement about the $abc$ conjecture, about Parshin's
conjecture, or about the truth of Theorems~1, A or~B of~\cite{So26} as statements.
It shows that the pencil step of a PMB argument cannot produce a positive margin
under the hypotheses stated below, and that the proof of~\cite[Theorem~1]{So26}
given in version~1 is incomplete at Proposition~9\,vi). All references
to~\cite{So26} concern version~1; a later version may differ. The intersection
numbers of~\cite[Prop.~8 and Prop.~9\,ii), iii)]{So26} that we use agree with our
own computation.

\section{Setting and the Margin of a Pencil}\label{sec:setting}

\subsection{Conventions}
In Sections~\ref{sec:setting}--\ref{sec:nef} we work over an algebraically closed
field of characteristic~$0$. Let $\phi\colon X\to Y$ be a surjective morphism
from a smooth projective surface onto a smooth projective curve whose fibers are
all integral. Let $F$ be the class of a fiber, $\eta$ the generic point of~$Y$,
and $\omega=\omega_{X/Y}$, and assume $\omega^2>0$. Smooth fibrations, in
particular Kodaira fibrations, satisfy these hypotheses. A divisor is
\emph{vertical} if its components lie in fibers and \emph{horizontal} if no
component does. Because the fibers are integral, every effective vertical divisor
is of the form $V=\sum_y n_yX_y$ with $n_y\ge 0$; we set $\deg V=\sum_y n_y$, so
that $V\equiv(\deg V)F$, $V^2=0$, and $V\cdot F=0$.

For a line bundle $L$ on $X$, set $m=L\cdot F\ge 1$ and $E=\phi_*L$. As a
torsion-free sheaf on the smooth curve $Y$, $E$ is locally free of rank
$r=h^0(X_\eta,L_\eta)$; we write $\mu(E)=\deg E/r$. Note that $L\cdot V=m\deg V$.
We will often assume:
\begin{itemize}[leftmargin=3.5em]
\item[\bpf] The linear system $|L_\eta|$ on the generic fiber has no base point.
\end{itemize}

\begin{lemma}\label{lem:image}
Let $0\ne E'\subseteq E$ be a subsheaf. The image of the composite
$\phi^*E'\to\phi^*E\to L$ is $L(-B'-V')\otimes\cI_{Z'}$ with $B'\ge 0$
horizontal, $V'\ge 0$ vertical, and $Z'$ a subscheme of finite length. The
restriction of $L(-B')$ to $X_\eta$ has a nonzero section. Under~\bpf{} one has
$B'=0$ for $E'=E$.
\end{lemma}

\begin{proof}
On a smooth surface, every nonzero subsheaf of a line bundle $L$ is of the form
$L(-D)\otimes\cI_{Z'}$ with $D\ge 0$ and $Z'$ of finite length; decompose
$D=B'+V'$ into its horizontal and vertical parts. Over $\eta$ the map is the
evaluation of the sections in $E'_\eta\subseteq E_\eta=H^0(X_\eta,L_\eta)$, so
its image is the subsheaf generated by these sections; its divisorial base locus
is $B'|_{X_\eta}$, and any nonzero section in $E'_\eta$ defines a section of
$L(-B')|_{X_\eta}$. For $E'=E$ this base locus is empty under~\bpf.
\end{proof}

\begin{definition}\label{def:margin}
Set $r_1=-L^2/\omega^2$ and $q=-\deg E/\omega^2$. The \emph{margin} of $L$ is
\[
M_r(L)=r_1-\frac{2m}{r}\,q=\frac{-L^2+\frac{2m}{r}\deg E}{\omega^2}.
\]
For $r=2$ we write $M(L)=M_2(L)=r_1-mq$, in accordance with~\eqref{eq:margin-intro}.
\end{definition}

\begin{lemma}\label{lem:square}
One has $M_r(L)=-(L-\mu(E)F)^2/\omega^2$ and $M_r(L+dF)=M_r(L)$ for every
integer $d$.
\end{lemma}

\begin{proof}
We have $(L-\mu F)^2=L^2-2m\mu=L^2-\frac{2m}{r}\deg E$. Twisting by $dF$ changes
$L^2$ by $2md$ and, by the projection formula, $\deg E$ by $rd$.
\end{proof}

\subsection{The Pencil Chain}\label{subsec:chain}
Now let $\phi\colon X\to Y$ be defined over a number field~$K$ with geometrically
integral fibers; the conventions above apply to $X_{\overline K}\to
Y_{\overline K}$, which leaves intersection numbers and degrees unchanged. For a
finite extension $F/K$ of degree $d$ and a point $P\in Y(F)$, let $X_P$ be the
fiber, equip $\omega$ and $L$ with adelic metrics as in~\cite[\S1.3]{So26}, and
let $H_P=H^0(X_P,L_P)$ carry the sup-norms. For $r=2$, the proof
of~\cite[Theorem~1]{So26} in~\cite[\S6.3]{So26} combines three steps:
\begin{enumerate}[label=(\roman*)]
\item \emph{Pencil inequality}~\cite[Prop.~3, (28), (29)]{So26}: if $s_1,s_2\in
H_P$ have no common zero, then $-\Lb_P^{\,2}\le m\,(h(s_1)+h(s_2))+O(d)$.
\item \emph{Minkowski}~\cite[(2), (30)]{So26}: there are linearly independent
$s_1,s_2\in H_P$ with $h(s_1)+h(s_2)\le-\degh(H_P)+\delta(F)+O(d)$.
\item \emph{Transfer}~\cite[(32), Cor.~1, (33), (35)]{So26}: for $P$ outside a
finite set and every $\varepsilon>0$,
$\degh(H_P)=-q\,\ob_P^{\,2}+O(\varepsilon\,\ob_P^{\,2})+O_\varepsilon(d)$ and
$\Lb_P^{\,2}=-r_1\,\ob_P^{\,2}+O(\varepsilon\,\ob_P^{\,2})+O_\varepsilon(d)$.
\end{enumerate}
Together they give
\begin{equation}\label{eq:chain}
\bigl(M(L)-O(\varepsilon)\bigr)\,\ob_P^{\,2}\le m\,\delta(F)+O_\varepsilon(d).
\end{equation}
For general $r$, two sections realizing the first two successive minima of $H_P$
satisfy $\lambda_1+\lambda_2\le\frac{2}{r}(\lambda_1+\dots+\lambda_r)$, and
Minkowski's second theorem bounds the sum by $-\degh(H_P)$ up to a multiple of
$\delta(F)$ and $O(d)$. If these two sections have no common zero, the same steps
yield~\eqref{eq:chain} with $M_r(L)$ in place of $M(L)$ and a modified constant
in front of $\delta(F)$.

\begin{remark}\label{rem:chain}
Inequality~\eqref{eq:chain} bounds $\ob_P^{\,2}$ from above if and only if
$M_r(L)>0$; for $M_r(L)\le 0$ it is vacuous. In~\cite{So26} one has $r=2$, $m=6$,
$L^2=-20k$, and $\omega^2=12k$, so $r_1=5/3$; Proposition~9\,vi) gives $q\le 1/4$,
hence $M(L)\ge 1/6$ and the constant $36=6/(1/6)$ of~\cite[Theorem~1]{So26}. The
$\varepsilon$- and $O(d)$-terms do not affect the sign. All results below are
statements about the classical invariants $L^2$, $\deg\phi_*L$, and $\omega^2$;
in this sense they are independent of the arithmetic details of the chain.
\end{remark}

\section{Rank Two: Xiao's Lemma for Pencils}\label{sec:rank2}

\begin{theorem}\label{thm:rank2}
In the setting of Section~\ref{sec:setting}, assume $r=2$ and~\bpf, and let
$L(-V)\otimes\cI_Z$ be the image of the evaluation map
$\phi^*E\to L$ (Lemma~\ref{lem:image}). Then
\begin{equation}\label{eq:rank2}
-L^2+m\deg E=-\ell(Z)-2m\deg V\le 0 .
\end{equation}
In particular, $\deg\phi_*L\le L^2/m$, with equality if and only if
$\phi^*\phi_*L\to L$ is surjective.
\end{theorem}

\begin{proof}
Set $A=L(-V)$, so that the evaluation map induces a surjection $\phi^*E\to
A\otimes\cI_Z$, and let $K$ be its kernel.

\emph{Step 1: $K$ is a line bundle.} At every point, the ideal $\cI_Z$ is a
torsion-free module over a two-dimensional regular local ring, hence of depth at
least one and, by the Auslander--Buchsbaum formula, of projective dimension at
most one. The kernel of a surjection from a locally free sheaf onto a sheaf of
projective dimension at most one is locally free; thus $K$ has rank one.

\emph{Step 2: Chern classes of the outer terms.} Since $\phi^*E$ is pulled back
from the curve $Y$, we have $c_1(\phi^*E)\equiv(\deg E)F$ and
$c_2(\phi^*E)=\phi^*c_2(E)=0$. From $0\to\cI_Z\to\OO_X\to\OO_Z\to 0$ and
$\ch(\OO_Z)=\ell(Z)[\mathrm{pt}]$ one obtains $c_1(\cI_Z)=0$ and
$c_2(\cI_Z)=\ell(Z)$; tensoring the rank-one sheaf $\cI_Z$ by $A$ gives
$c_1(A\otimes\cI_Z)=A$ and $c_2(A\otimes\cI_Z)=\ell(Z)$.

\emph{Step 3: Whitney formula.} We have $c_1(K)\equiv(\deg E)F-A$, and
\[
0=c_2(\phi^*E)=c_1(K)\cdot A+\ell(Z)=m\deg E-A^2+\ell(Z),
\]
using $F\cdot A=m$ (since $F\cdot V=0$).

\emph{Step 4.} We have $A^2=(L-V)^2=L^2-2m\deg V$, since $V^2=0$ and $L\cdot V=m\deg
V$. Substituting gives~\eqref{eq:rank2}. Equality $\deg E=L^2/m$ holds if and
only if $Z=\emptyset$ and $V=0$, that is, if and only if the evaluation map is
surjective.
\end{proof}

\begin{corollary}\label{cor:rank2}
Under the hypotheses of Theorem~\ref{thm:rank2},
\[
M(L)=-\frac{\ell(Z)+2m\deg V}{\omega^2}=-\frac{(L-\mu(E)F)^2}{\omega^2}\le 0 .
\]
The pencil chain~\eqref{eq:chain} is therefore vacuous for every line bundle of
rank two with base-point-free generic linear system on every fibration with
integral fibers.
\end{corollary}

\begin{remark}\label{rem:rank2}
(a) Theorem~\ref{thm:rank2} requires no stability hypothesis on $\phi_*L$; it is
the exact rank-two form of Xiao's method~\cite{Xi87}. Equivalently,
$(L-\mu(E)F)^2=\ell(Z)+2m\deg V\ge 0$.

(b) For $r=2$, no cleverer choice of the two sections helps. Minkowski's second
theorem determines $\lambda_1+\lambda_2$ up to terms of order $\delta(F)+d$, so
together with step~(iii) every pair of linearly independent sections satisfies
$h(s_1)+h(s_2)\ge q\,\ob_P^{\,2}-O(\varepsilon\,\ob_P^{\,2}+\delta(F)+d)$. The
best coefficient the chain can achieve is $r_1-mq=M(L)\le 0$.

(c) Twisting by vertical divisors does not help either: $M(L+dF)=M(L)$ by
Lemma~\ref{lem:square}.
\end{remark}

\section{Rank \texorpdfstring{$r$}{r}: A Margin Identity}\label{sec:rankr}

For a locally free sheaf $G$ of rank $n$ on $X$, let
$\Delta(G)=2n\,c_2(G)-(n-1)\,c_1(G)^2=c_1(G)^2-2n\ch_2(G)$ denote its Bogomolov
discriminant; $\Delta(G)=0$ for line bundles.

\begin{theorem}\label{thm:rankr}
In the setting of Section~\ref{sec:setting}, assume~\bpf{} and $r\ge 2$. Let
$L(-V)\otimes\cI_Z$ be the image of $\phi^*E\to L$ and $K$ its kernel. Then $K$ is
locally free of rank $r-1$, and
\begin{equation}\label{eq:rankr}
\Delta(K)=r\,L^2-2m\deg E-2rm\deg V-2(r-1)\ell(Z),
\end{equation}
or equivalently
\[
M_r(L)=-\frac{\Delta(K)+2rm\deg V+2(r-1)\ell(Z)}{r\,\omega^2}.
\]
\end{theorem}

\begin{proof}
Local freeness follows as in Step~1 of Theorem~\ref{thm:rank2}. Set $A=L(-V)$.
From $c(\phi^*E)=c(K)\,c(A\otimes\cI_Z)$ and Step~2 above, we have
$c_1(K)\equiv(\deg E)F-A$ and $0=c_2(K)+c_1(K)\cdot A+\ell(Z)$, whence
\[
c_1(K)^2=A^2-2m\deg E,\qquad c_2(K)=A^2-m\deg E-\ell(Z).
\]
Therefore
$\Delta(K)=2(r-1)c_2(K)-(r-2)c_1(K)^2=r\,A^2-2m\deg E-2(r-1)\ell(Z)$, and
substituting $A^2=L^2-2m\deg V$ gives~\eqref{eq:rankr}. Dividing by $-r\,\omega^2$
yields the expression for $M_r(L)$. For $r=2$ the kernel is a line bundle,
$\Delta(K)=0$, and~\eqref{eq:rankr} reduces to~\eqref{eq:rank2}.
\end{proof}

\begin{corollary}\label{cor:bogomolov}
If, in addition, $K$ is $\mu$-semistable with respect to some ample class on $X$,
then Bogomolov's inequality $\Delta(K)\ge 0$~\cite{Bo78} (see
also~\cite[\S3.4]{HL10}) gives $M_r(L)\le 0$.
\end{corollary}

\begin{remark}\label{rem:rankr}
Semistability of the syzygy kernel is not automatic, and for $r\ge 3$ the
estimate $L^2\ge\frac{2m}{r}\deg\phi_*L$ is a slope inequality that is derived in
the literature under stability hypotheses (compare~\cite{Xi87,CH88}).
Section~\ref{sec:nef} replaces the hypothesis on $K$ by semistability of
$\phi_*L$ on the base curve and also treats the unstable case.
\end{remark}

\section{The Nef Lemma and the Harder--Narasimhan Dichotomy}\label{sec:nef}

Let $0=E_0\subsetneq E_1\subsetneq\dots\subsetneq E_n=E$ be the
Harder--Narasimhan filtration of $E=\phi_*L$~\cite{HN75}, with slopes
$\mu_j=\mu(E_j/E_{j-1})$, so that $\mu_1>\dots>\mu_n$ and $\mu_j$ is the minimal
slope of $E_j$. By Lemma~\ref{lem:image}, the image of $\phi^*E_j\to L$ is
$L(-B_j-V_j)\otimes\cI_{Z_j}$; if~\bpf{} holds, then $B_n=0$, $V_n=V$, and
$Z_n=Z$. Set
\[
N_j=L-B_j-V_j-\mu_jF\in\Pic(X)\otimes\mathbb{Q}.
\]

\begin{lemma}[nef lemma]\label{lem:nef}
For every $j$, the $\mathbb{Q}$-divisor $N_j$ is nef. Hypothesis~\bpf{} is not
needed.
\end{lemma}

\begin{proof}
Let $C\subset X$ be an irreducible curve. If $C$ is vertical, then $C$ is a fiber
and $N_j\cdot C=(L-B_j)\cdot F\ge 0$, because $L(-B_j)|_{X_\eta}$ has a nonzero
section (Lemma~\ref{lem:image}). Now let $C$ be horizontal, $\nu\colon\widetilde C\to
C$ its normalization, and $f=\phi\circ\nu\colon\widetilde C\to Y$, which is a
finite morphism of degree $d_C=C\cdot F\ge 1$. The map $\phi^*E_j\to
L(-B_j-V_j)\otimes\cI_{Z_j}$ is surjective onto $L(-B_j-V_j)$ outside the finite
set $Z_j$, in particular at the generic point of $C$. Pulling back to
$\widetilde C$ gives a nonzero homomorphism $f^*E_j\to\nu^*L(-B_j-V_j)$ whose
image $Q$ is an invertible quotient of $f^*E_j$ contained in $\nu^*L(-B_j-V_j)$.
Hence
\[
(L-B_j-V_j)\cdot C=\deg\nu^*L(-B_j-V_j)\ge\deg Q\ge\mu_{\min}(f^*E_j)=d_C\,\mu_j .
\]
The last equality holds because in characteristic~$0$ the pullback of the
Harder--Narasimhan filtration under a finite morphism of smooth curves coincides
with the Harder--Narasimhan filtration of the pullback, with slopes multiplied by
the degree (see~\cite[\S3.2]{HL10}). Since $F\cdot C=d_C$, this gives
$N_j\cdot C\ge 0$.
\end{proof}

\begin{corollary}[semistable case]\label{cor:ss}
Assume~\bpf{} and that $E=\phi_*L$ is semistable. Then $N=L-V-\mu(E)F$ is nef, and
\[
M_r(L)=-\frac{N^2+2m\deg V}{\omega^2}\le 0 .
\]
No hypothesis on the stability of the fibers, of their linear systems, or of the
syzygy kernel is needed.
\end{corollary}

\begin{proof}
Here $n=1$, $E_1=E$, $B_1=0$, $V_1=V$, and $\mu_1=\mu(E)$, so $N=N_1$ is nef by
Lemma~\ref{lem:nef} and hence $N^2\ge 0$. Since $L-\mu(E)F=N+V$,
$N\cdot V=m\deg V$, and $V^2=0$, we obtain $(L-\mu(E)F)^2=N^2+2m\deg V$; the
claim follows from Lemma~\ref{lem:square}.
\end{proof}

\begin{proposition}[unstable case]\label{prop:unstable}
Assume~\bpf{} and that $E=\phi_*L$ is unstable, $n\ge 2$.
\begin{enumerate}[label=(\alph*)]
\item If $B_1=0$, then
$M_r(L)\le-2m\bigl(\mu_1-\mu(E)+\deg V_1\bigr)/\omega^2<0$.
\item For all indices $a,b$ with $B_b=0$,
\[
L^2\ge m(\mu_a+\mu_b)+m\deg V_b ,\qquad\text{that is}\qquad
\frac{-L^2+m(\mu_a+\mu_b)}{\omega^2}\le-\frac{m\deg V_b}{\omega^2}\le 0 .
\]
Such an index $b$ always exists, since $B_n=0$.
\end{enumerate}
\end{proposition}

\begin{proof}
(a) We have $L-\mu(E)F=N_1+V_1+(\mu_1-\mu(E))F$. Expanding with
$N_1\cdot V_1=m\deg V_1$, $N_1\cdot F=m$ (as $B_1=0$), and
$V_1^2=V_1\cdot F=F^2=0$ gives
$(L-\mu(E)F)^2=N_1^2+2m\deg V_1+2m(\mu_1-\mu(E))$. Since $N_1^2\ge 0$ and
$\mu_1>\mu(E)$ for unstable $E$, the claim follows from Lemma~\ref{lem:square}.

(b) The divisors $N_b$ and $N_a$ are nef and $B_a$, $V_a$ are effective, so
$N_b\cdot(N_a+B_a+V_a)\ge 0$. Here $N_a+B_a+V_a=L-\mu_aF$ and $N_b=L-V_b-\mu_bF$,
and the intersection product equals $L^2-m\mu_a-m\mu_b-m\deg V_b$.
\end{proof}

\begin{remark}\label{rem:unstable}
In case (a), the margin that the chain of Section~\ref{subsec:chain} produces with
Minkowski's theorem on the full lattice is strictly negative. The coefficients in
(b) are those that the chain would produce if the heights of the two sections were
governed by the slopes $\mu_a$ and $\mu_b$; for instance, Corollary~1
of~\cite{So26}, applied to the sublattice coming from $E_1$, gives a section $s_1$
with
$h(s_1)\le(-\mu_1/\omega^2+\varepsilon)\,\ob_P^{\,2}+c\,\delta(F)+O_\varepsilon(d)$
for a constant $c$. Thus every coefficient that the chain can obtain from
geometric subbundles in this way is nonpositive. The condition $B_b=0$ in (b)
reflects the pencil step: every section coming from $E_j$ vanishes along
$B_j\cap X_P$, so two sections taken from the same step $E_a$ with $B_a\ne 0$
always have a common zero on $X_P$. If $B_a$ and $B_b$ merely have no common
component, their restrictions to $X_P$ are disjoint for all but finitely many $P$,
and the same nonpositivity follows from Proposition~\ref{prop:cross} with
$L_1=L_2=L$.
\end{remark}

\begin{problem}[arithmetic anomaly, $r\ge 3$]\label{prob:anomaly}
Let $r\ge 3$ and assume that $\phi_*L$ is unstable. We do not know whether, for
points $P$ of large height, an upper bound for $h(s_1)+h(s_2)$ over pairs of
sections without common zero can be obtained that lies below what the geometric
slopes $\mu_a$ predict. Excluding this would require a comparison between the
arithmetic Harder--Narasimhan polygon of the adelic lattice $H_P$ and the
geometric Harder--Narasimhan polygon of $\phi_*L$, scaled by
$\ob_P^{\,2}/\omega^2$. We leave this question open. For $r=2$ it does not arise
(Remark~\ref{rem:rank2}\,(b)), and for semistable $\phi_*L$ it is answered by
Corollary~\ref{cor:ss}.
\end{problem}

\section{Application to Soul\'e, arXiv:2611.05217v1}\label{sec:soule}

\subsection{The Construction}
We recall the Kodaira fibration of~\cite[\S\S3, 5]{So26}, which
follows~\cite{Ko67,MD85}. Let $Y_0$ be a curve of genus~$2$ and $\chi\colon
Y_1\to Y_0$ a nontrivial \'etale double cover with covering involution $\nu$;
$Y_1$ has genus~$3$ and is hyperelliptic~\cite[Prop.~4\,i)]{So26}. Let $\iota$ be
its hyperelliptic involution and $H$ the hyperelliptic class; then
$\omega_{Y_1}\cong\OO(2H)\cong\OO(4W)$ for a Weierstrass point $W$. Restricting
the relation $\OO(\Gamma-2\Delta_1)=(1\times\psi_1)^*\mathcal{P}+\pi_2^*L$
of~\cite[\S5.1]{So26} to $Y_1\times\{x\}$ shows that $\psi_1\colon Y_1\to J_1$ is
given by $x\mapsto[\nu x-x]$. The curve $Y$ is a connected component of the fiber
product $Z$ of $\psi_1$ with multiplication by $2$ on $J_1$, and
$\psi\colon Y\to Y_1$ is finite \'etale of degree~$k$.

On $S=Y_1\times Y$ with projections $\mathrm{pr}_1$ and $p\colon S\to Y$, let
$D_1$ and $D_2$ be the graphs of $\psi$ and $\nu\circ\psi$; they are disjoint
because $\nu$ has no fixed point. Let $\cL\in\Pic(S)$ with
$2\cL=\OO_S(D_1+D_2)$~\cite{MD85}, let $\rho\colon X\to S$ be the double cover
branched along $D_1+D_2$, and set $\phi=p\circ\rho$; this morphism is smooth with
fibers of genus~$6$. Let $E_i=\rho^{-1}(D_i)_{\mathrm{red}}$, so that
$\rho^*D_i=2E_i$, and set $B_i=\OO_X(E_i)$. For $y\in Y$ set $P_1=\psi(y)$,
$P_2=\nu P_1$, and $\lambda_y=\cL|_{Y_1\times\{y\}}$, a square root of
$\OO(P_1+P_2)$. The point $P_1$ is \emph{exceptional} if $h^0(P_1+P_2)=2$;
equivalently, $P_2=\iota P_1$~\cite[Prop.~5\,iii)]{So26}. There are four
exceptional points, and $\Psi=\psi^{-1}(\text{exceptional points})$ has $4k$
elements.

Let $G=\phi_1^*\OO(W)$ with $\phi_1=\mathrm{pr}_1\circ\rho$. The lattice spanned
by $G,B_1,B_2,F$ has the intersection numbers in Table~\ref{tab:lattice}; they
are those of~\cite[Prop.~8]{So26} together with $B_i\cdot F=1$ ($E_i$ is a
section) and $G\cdot F=2$ ($\phi_1$ has degree two on fibers). The
Riemann--Hurwitz formula gives $\omega=4G+B_1+B_2$ and $\omega^2=12k$. Soul\'e's
auxiliary bundle is
\[
L=\omega-2B_1-2B_2=4G-B_1-B_2,\qquad m=6,\qquad L^2=-20k,
\]
and $h^0(L_y)=2$ for general $y$~\cite[Prop.~9\,ii), iii), v)]{So26}.

\begin{table}[ht]
\centering
\begin{tabular}{@{}lcccc@{}}
\toprule
 & $G$ & $B_1$ & $B_2$ & $F$\\
\midrule
$G$ & $0$ & $k$ & $k$ & $2$\\
$B_1$ & $k$ & $-2k$ & $0$ & $1$\\
$B_2$ & $k$ & $0$ & $-2k$ & $1$\\
$F$ & $2$ & $1$ & $1$ & $0$\\
\bottomrule
\end{tabular}
\caption{Intersection numbers on the Kodaira fibration of~\cite{So26}.}
\label{tab:lattice}
\end{table}

\subsection{The Count}

\begin{lemma}[$k=4$]\label{lem:k4}
Let $G_2\subseteq J_1[2]$ be the monodromy group of the $J_1[2]$-torsor $Z\to Y_1$.
Then $G_2=(\nu_*-1)\,J_1[2]$ and $|G_2|=4$. Consequently, $k=4$ for every choice
of $Y_0$, of $\chi$, and of the component $Y$.
\end{lemma}

\begin{proof}
The torsor is the pullback of $[2]\colon J_1\to J_1$ along $\psi_1$, so its
monodromy is $\pi_1(Y_1)\to H_1(Y_1,\mathbb{Z})\xrightarrow{(\psi_1)_*}
H_1(J_1,\mathbb{Z})\to J_1[2]$ with $(\psi_1)_*=\nu_*-1$, and every connected
component has degree $|G_2|$ over $Y_1$. The order can be computed over
$\mathbb{C}$. Since the mapping class group acts transitively on the nonzero
classes in $H^1(Y_0,\mathbb{Z}/2)$, we may take
$\pi_1(Y_0)=\langle a_1,b_1,a_2,b_2\mid[a_1,b_1][a_2,b_2]\rangle$ and assume that
$\chi$ is given by the character with $a_1\mapsto 1$ and $b_1,a_2,b_2\mapsto 0$.
With the Schreier transversal $\{1,a_1\}$, the subgroup $\pi_1(Y_1)$ is generated
by $s=a_1^2$, $b_1$, $a_2$, $b_2$ and their conjugates $b_1'$, $a_2'$, $b_2'$ by
$a_1$. The defining relator and its conjugate both map to $b_1'=b_1$ in the
abelianization, so $H_1(Y_1,\mathbb{Z})$ is free on $s,b_1,a_2,b_2,a_2',b_2'$.
The involution $\nu$ corresponds to conjugation by $a_1$; on homology it fixes
$s$ and $b_1$ and exchanges $a_2\leftrightarrow a_2'$ and
$b_2\leftrightarrow b_2'$. The image of $\nu_*-1$ is spanned by $a_2'-a_2$ and
$b_2'-b_2$, which extend to a basis, so its reduction modulo~$2$ has order~$4$.
\end{proof}

\begin{lemma}[effective square roots]\label{lem:keff}
Let $P_1$ be exceptional. The classes $\lambda_y$, $y\in\psi^{-1}(P_1)$, form a
coset of $G_2$ in the set of square roots of $H$. Every coset of $G_2$ in this set
contains at most two effective classes.
\end{lemma}

\begin{proof}
For $(x,z)\in Z$, the class $x+z\in\Pic^1(Y_1)$ is a square root of $x+\nu x$,
since $2(x+z)=2x+\nu x-x$. The map $y\mapsto\lambda_y-(\psi(y)+\psi_2(y))$, with
$\psi_2\colon Y\to J_1$ denoting the second projection, takes values in $J_1[2]$
and is therefore constant on the connected curve $Y$. Over $P_1$ the points of $Y$
form a $G_2$-orbit in the fiber of $Z$, so the classes $\lambda_y$ form a coset of
$G_2$. An effective square root of $H$ has the form $\OO(Q)$ with $2Q\in|H|$, and
thus corresponds to a Weierstrass point. Let $W\ne W'$ be effective in the same
coset. Then $W'-W=(\nu-1)g$ for some $g\in J_1[2]$, hence
$(1+\nu)(W'-W)=(\nu^2-1)g=0$ and $W'+\nu W'\sim W+\nu W$. If these divisors were
distinct, $W+\nu W$ would move in a pencil and hence lie in $|H|$, giving
$\nu W=\iota W=W$, which is impossible because $\nu$ is fixed-point-free. Hence
$W'=\nu W$, and a coset contains at most the two effective roots $W,\nu W$.
\end{proof}

\begin{lemma}[torsion of length one]\label{lem:length}
The torsion of $R^1p_*\cL$ is supported precisely at the points $y$ with
$h^0(\lambda_y)\ne 0$, and has length one at each of them.
\end{lemma}

\begin{proof}
Generically $h^0(\lambda_y)=0$~\cite[Prop.~5\,ii)]{So26}, so $p_*\cL=0$. Let $\pi$
be a uniformizer at $y$. The exact sequence $0\to\cL\xrightarrow{\pi}\cL\to\lambda_y\to
0$ gives $R^1p_*\cL[\pi]\cong H^0(\lambda_y)$, so the torsion $T_y$ is nonzero
precisely when $h^0(\lambda_y)\ne 0$, and it is cyclic because
$h^0(\lambda_y)\le 1$. Likewise, $R^1p_*\cL[\pi^2]$ is isomorphic to the space of
sections of $\cL$ over $Y_1\times\operatorname{Spec}\OO_{Y,y}/\pi^2$, so
$\ell(T_y)\ge 2$ if and only if the section $s_W$ of $\lambda_y=\OO(W)$ extends to
first order. This holds if and only if $v\cup s_W=0$ in $H^1(\OO(W))$, where
$v\in H^1(\OO_{Y_1})$ is the tangent vector of $y\mapsto\lambda_y$; that is, if
and only if $v$ lies on the line $d\mathrm{AJ}_W(T_WY_1)$. Under $H^1(\OO_{Y_1})
\cong H^0(\omega_{Y_1})^\vee$, this line is spanned by the evaluation at $W$.
Since $2\lambda_y=\OO(\psi(y)+\nu\psi(y))$ and $\psi$ is \'etale, $2v$ is the
functional $\alpha\mapsto\alpha_{P_1}(t)+\alpha_{P_2}(d\nu\,t)$ for a nonzero
tangent vector $t$ at $P_1$. At an exceptional point $P_2=\iota P_1$, since
$\iota$ commutes with $\nu$, the involution $\nu\iota$ fixes $P_2$ and has
derivative $-1$ there, so $d\nu=-d\iota$ at $P_1$; as $\iota^*\alpha=-\alpha$ for
all $\alpha\in
H^0(\omega_{Y_1})$, we obtain $2v=2\,\mathrm{ev}_{P_1,t}\ne 0$. The evaluations at
$P_1$ and $W$ are proportional only if the canonical map identifies $P_1$ and $W$,
that is, if $P_1\in\{W,\iota W\}=\{W\}$. However, $P_1=W$ would give $W\sim P_2$
and hence $P_2=P_1$, a contradiction. Thus $\ell(T_y)=1$.
\end{proof}

\begin{proposition}[degree formula]\label{prop:degree}
Let $T$ be the number of points $y\in Y$ with $h^0(\lambda_y)\ne 0$; all of them
lie in $\Psi$. Then
\[
\deg\phi_*L=-7k+T .
\]
\end{proposition}

\begin{proof}
Set $K_1=\mathrm{pr}_1^*\omega_{Y_1}=\omega_{S/Y}$. The canonical bundle formula
for double covers~\cite[I.17]{BHPV04} gives $\omega_{X/Y}=\rho^*(K_1+\cL)$, and
$B=2B_1+2B_2=\rho^*(D_1+D_2)=\rho^*(2\cL)$. Hence $L=\rho^*(K_1-\cL)$, and by
$\rho_*\OO_X=\OO_S\oplus\cL^{-1}$~\cite[(13)]{So26} we have
\[
\phi_*L=p_*(K_1-\cL)\oplus p_*(K_1-2\cL).
\]
\emph{Second summand.} Since $x+\iota x\in|H|$ for every $x$, the seesaw principle
gives $\OO_S(D_i+\iota D_i)\cong\mathrm{pr}_1^*\OO(H)\otimes p^*M_i$ with
$M_i\in\Pic(Y)$, where $\iota D_i=(\iota\times 1)(D_i)$. Intersecting with
$\{x\}\times Y$, which meets each of $D_i$ and $\iota D_i$ transversally in $k$
points and on which $\mathrm{pr}_1^*\OO(H)$ is trivial, gives $\deg M_i=2k$. As
$K_1=\mathrm{pr}_1^*\OO(2H)$, we have
\[
K_1-2\cL=K_1-D_1-D_2\cong\OO_S(\iota D_1+\iota D_2)\otimes p^*(-M_1-M_2).
\]
The automorphism $\iota\times 1$ of $S$ over $Y$ and~\cite[(20)]{So26} give
$p_*\OO_S(\iota D_1+\iota D_2)\cong p_*\OO_S(D_1+D_2)\cong\OO_Y$, hence
$\deg p_*(K_1-2\cL)=-4k$.

\emph{First summand.} Relative duality gives
$p_*(K_1-\cL)\cong\mathcal{H}om(R^1p_*\cL,\OO_Y)$, so
$\deg p_*(K_1-\cL)=-(\deg R^1p_*\cL-T)$ by Lemma~\ref{lem:length}. Since
$p_*\cL=0$, Grothendieck--Riemann--Roch for $p$ gives
$\deg R^1p_*\cL=-\bigl(\tfrac12\cL^2-\tfrac12K_1\cdot\cL\bigr)$, and
$4\cL^2=(D_1+D_2)^2=-8k$ together with $2K_1\cdot\cL=K_1\cdot(D_1+D_2)=8k$ yields
$\deg
R^1p_*\cL=3k$. Hence $\deg p_*(K_1-\cL)=T-3k$.

Adding the two degrees gives $\deg\phi_*L=T-3k-4k$. By~\cite[Prop.~5\,ii)]{So26},
$h^0(\lambda_y)\ne 0$ forces $P_1$ to be exceptional, so the points counted by $T$
lie in $\Psi$.
\end{proof}

\begin{theorem}\label{thm:soule-count}
In the Kodaira fibration of~\cite{So26}, with $L=\omega-2B_1-2B_2$,
\[
\deg\phi_*L=-7k+4k_{\mathrm{eff}},\qquad k=4,\qquad 0\le k_{\mathrm{eff}}\le 2,
\]
where $4k_{\mathrm{eff}}=T$. In particular, $\deg\phi_*L\le-20<-12=-3k$, so
Proposition~9\,vi) of arXiv:2611.05217v1 is false. Moreover,
$q=-\deg\phi_*L/\omega^2\ge 5/12$, and the leading coefficient $5/3-6q$
of~\cite[(36)]{So26} satisfies $5/3-6q\le-5/6$.
\end{theorem}

\begin{proof}
Proposition~\ref{prop:degree} gives $\deg\phi_*L=-7k+T$. Lemma~\ref{lem:k4} gives
$k=4$, and Lemma~\ref{lem:keff} shows that each of the four fibers of $\psi$ over
an exceptional point contributes at most two points to $T$, so $T\le 8$. Hence
$\deg\phi_*L\le-28+8=-20$ and $q\ge 20/48=5/12$.
\end{proof}

\begin{theorem}[structural form]\label{thm:soule-structural}
In the setting of~\cite{So26}, the linear system $|L_\eta|$ is base-point-free for
every choice of the component $Y$ and of $\cL$. Consequently,
$\deg\phi_*L\le L^2/6=-10k/3<-3k$ and $M(L)\le 0$; in particular, no choice of the
component $Y$ or of $\cL$ satisfies Proposition~9\,vi).
\end{theorem}

\begin{proof}
\emph{Step 1: fibers outside $\Psi$.} Let $y\notin\Psi$, so that $P_2\ne\iota P_1$.
Then $h^0(\lambda_y)=0$ by~\cite[Prop.~5\,ii)]{So26}, and Riemann--Roch gives
$h^0(\omega_{Y_1}-\lambda_y)=1$ and $h^0(\omega_{Y_1}-2\lambda_y)=h^0(P_1+P_2)=1$.
As $\omega_{Y_1}\cong\OO(2H)$ and $P_i+\iota P_i\in|H|$, the unique effective
divisor in $|\omega_{Y_1}-2\lambda_y|$ is $\iota P_1+\iota P_2$. By the
decomposition of $\phi_*L$ in the proof of Proposition~\ref{prop:degree},
$H^0(X_y,L_y)$ is spanned by $\rho^*\sigma_y$ and $t_y\,\rho^*\tau_y$, where
$\sigma_y$ spans $H^0(\omega_{Y_1}-\lambda_y)$, $\tau_y$ has divisor
$\iota P_1+\iota P_2$, and $t_y$ is the section of $\rho^*\lambda_y$ that vanishes
at the two ramification points. Hence $|L_y|$ has a base point if and only if
$\sigma_y$ vanishes at one of $P_1$, $P_2$, $\iota P_1$, $\iota P_2$, and by
Riemann--Roch $\sigma_y$ vanishes at a point $Q$ if and only if
$h^0(\lambda_y+Q)\ne 0$.

\emph{Step 2: specialization.} On the open subset $Y^\circ\subseteq Y$ where
$h^0(\lambda_y)=0$, one has $h^0(\omega_{Y_1}-\lambda_y)=1$, so $p_*(K_1-\cL)$ is
a line bundle on $Y^\circ$ whose formation commutes with base change. Hence, for
each of the four maps $Q=\psi,\nu\psi,\iota\psi,\iota\nu\psi$ from $Y$ to $Y_1$,
the set of $y\in
Y^\circ$ with $\sigma_y(Q(y))=0$ is closed in $Y^\circ$. None of these sets is all
of $Y^\circ$: indeed, let $x_0$ be exceptional; by Lemmas~\ref{lem:k4}
and~\ref{lem:keff} there exists $y_0\in\psi^{-1}(x_0)$ with $h^0(\lambda_0)=0$,
where $\lambda_0=\lambda_{y_0}$. Since $\nu x_0=\iota x_0\ne x_0$, the four points
$Q(y_0)$ lie in $\{x_0,\iota
x_0\}$, neither of them is a Weierstrass point, and $2\lambda_0\sim x_0+\iota
x_0\sim H$. If $\lambda_0+q\sim a+b$ with $q\in\{x_0,\iota x_0\}$, then
$2a+2b\in|H+2q|$. As $q$ is not a Weierstrass point, $h^0(H+2q)=2$ and
$|H+2q|=|H|+2q$, so $2a+2b=c+\iota c+2q$ for some point $c$. All multiplicities on
the left-hand side are even, whereas for $c\ne\iota c$ either $c$ or $q$ has odd
multiplicity on the right-hand side. Hence $c=\iota c$ and $a+b=c+q$, so
$\lambda_0\sim c$ is effective, a contradiction. Thus $\sigma_{y_0}$ vanishes at
none of the points $Q(y_0)$, the four closed subsets are finite, and by Step~1 the
linear system $|L_y|$ is base-point-free for all but finitely many
$y\in Y\setminus\Psi$.

\emph{Step 3: the generic fiber.} For $y\in Y\setminus\Psi$ one has $h^0(L_y)=2$,
so $\phi_*L\otimes k(y)\to H^0(X_y,L_y)$ is bijective. The closure of a base point
of $|L_\eta|$ is a horizontal curve along which $\phi^*\phi_*L\to L$ is not
surjective; it meets every fiber $X_y$, and for $y\in Y\setminus\Psi$ it meets $X_y$
in a base point of $|L_y|$, contrary to Step~2. Hence $|L_\eta|$ is
base-point-free, and Theorem~\ref{thm:rank2} with $m=6$ and $L^2=-20k$ gives the
claimed bound.

\emph{Step 4: relation to~\cite{So26}.} In~\cite[\S6.3.3]{So26} the square root
$\cL'$ is chosen depending on the point $\bar y'\notin\Phi'$, and~\cite[Lemma~3]{So26}
shows that $L'_{P'}$ has no base point. Since $h^0(L'_{P'})=2$ equals the rank of
$\phi'_*L'$, the argument of Step~3 shows that this already forces $|L'_\eta|$ to be
base-point-free; thus every bundle to which~\cite[(28)]{So26} is applied falls under
Theorem~\ref{thm:rank2}. Steps~1--3 apply verbatim after the base change $a\colon
Y'\to Y$ of~\cite[\S6.3.1]{So26}, with $\psi\circ a$ in place of $\psi$. Let $J'$ be
the Jacobian of $Y'$. Since $\operatorname{Hom}(J_1,J')$ is torsion-free,
$\Pic(Y_1\times Y')[2]=\mathrm{pr}_1^*J_1[2]\oplus\mathrm{pr}_2^*J'[2]$, so on fibers
$\cL'$ restricts to the classes $\lambda_{a(y')}+\varepsilon$ for a fixed
$\varepsilon\in J_1[2]$. Over an exceptional point $x_0$, the classes
$\lambda_y+\varepsilon$, $y\in\psi^{-1}(x_0)$, form a coset of $G_2$, so by
Lemma~\ref{lem:keff} some $y_0\in\psi^{-1}(x_0)$ satisfies
$h^0(\lambda_{y_0}+\varepsilon)=0$; as $a$ is surjective, $y_0$ lifts to a point of
$(\psi\circ a)^{-1}(x_0)$, which is all that Step~2 requires. Lemma~\ref{lem:k4} is
used only for $\psi$, with $k=4$. Base-point-freeness therefore does not depend on
the pointwise choice.
\end{proof}

\begin{remark}[where the proof of Proposition~9 breaks down]\label{rem:gap}
Inequality~(24) of~\cite{So26} asserts that
$\deg(R^1\phi_*B)_{\mathrm{tor}}\ge 8k$. It rests on~\cite[Prop.~9\,iv)]{So26},
which states that $h^0(B_{\bar y})=3$ whenever $P_1$ is exceptional.
Its proof uses the sentence ``when $P_1$ is exceptional $h^0(\lambda)=1$'', with
a reference to parts ii) and iii) of~\cite[Prop.~5]{So26}. Proposition~5\,ii),
however, states that $h^0(\lambda)=0$ unless $\lambda$ is a Weierstrass point, in
which case $h^0(\lambda)=1$ and $h^0(P_1+P_2)=2$; it thus proves ``$\lambda$
effective $\Rightarrow$ $P_1$ exceptional'', the converse of the implication
needed. By Lemma~\ref{lem:keff}, the implication used fails at no fewer than
$4k-8=8$ of the $16$ points of $\Psi$; there $h^0(B_{\bar y})=2$ and
$h^0(L_{\bar y})=3$, so Proposition~9\,iv) and~v) fail at these points. Combined
with~\cite[(16), (18), (22)]{So26}, Proposition~\ref{prop:degree} gives
$\deg(R^1\phi_*B)_{\mathrm{tor}}=4k+T\le 24$ instead of $\ge 32$. Since a global
modification of $\cL$ by an element of $J_1[2]$ translates all cosets
simultaneously, no choice of $\cL$ alters the bound of Lemma~\ref{lem:keff}.
Theorem~1 of~\cite{So26} remains open as a statement; what fails is its proof via
Proposition~9\,vi), and by Corollary~\ref{cor:rank2} together with Step~4 of the
proof of Theorem~\ref{thm:soule-structural} the pencil step cannot be salvaged by
another rank-two auxiliary bundle $L'$ with base-point-free $|L'_\eta|$. If
$|L'_\eta|$ has a base point, then $\phi_*L'\otimes k(P)\to H^0(X_P,L'_P)$ is
bijective for $P$ outside a finite set, so by the argument of Step~3 of that proof
any two sections of $L'_P$ have a common zero for all such $P$, and the pencil
hypothesis of~\cite[(28)]{So26} cannot be satisfied.
\end{remark}

\section{Scope: Cross Terms, Single Sections, Source Bundles}\label{sec:scope}

This section records three further consequences of the nef lemma and of the
identity underlying Theorem~\ref{thm:rankr}. The formal statements are geometric
and proved in full; the remarks on their arithmetic meaning are interpretive.

\subsection{Cross Terms}
Proposition~3 of~\cite{So26} is stated for two line bundles $L_1,L_2$ of fiber
degrees $d_1,d_2$: $-\Lb_1\cdot\Lb_2\le d_2h(s_1)+d_1h(s_2)$. If $s_1$ comes from
the Harder--Narasimhan step $a$ of $\phi_*L_1$ and $s_2$ from step $b$ of
$\phi_*L_2$, with heights governed by the slopes $\mu_{1,a}$ and $\mu_{2,b}$, the
leading coefficient becomes
\[
M(L_1,L_2)=\frac{-L_1\cdot L_2+d_2\mu_{1,a}+d_1\mu_{2,b}}{\omega^2}
=-\frac{(L_1-\mu_{1,a}F)\cdot(L_2-\mu_{2,b}F)}{\omega^2}.
\]

\begin{proposition}\label{prop:cross}
If the horizontal base divisors $B_{1,a}$ and $B_{2,b}$ have no common component,
then $(L_1-\mu_{1,a}F)\cdot(L_2-\mu_{2,b}F)\ge 0$, and hence $M(L_1,L_2)\le 0$.
\end{proposition}

\begin{proof}
Write $L_1-\mu_{1,a}F=N_{1,a}+B_{1,a}+V_{1,a}$ and similarly for $L_2$, with
$N_{1,a},N_{2,b}$ nef by Lemma~\ref{lem:nef}. In the expansion, every term is
nonnegative: the intersection of two nef divisors, of a nef and an effective
divisor, $B_{1,a}\cdot B_{2,b}\ge 0$ (since they share no common component), of a
horizontal and a vertical divisor, and $V_{1,a}\cdot V_{2,b}=0$.
\end{proof}

\subsection{Single Sections and Jets along Curves}
For a single section $s$ and an irreducible horizontal curve $C$ not contained in
the divisor of $s$, comparing pointwise norms with sup-norms gives a bound of the
form $-\Lb\cdot\overline C\le(C\cdot F)\,h(s)$, whose geometric coefficient is
$-(L-\mu_jF)\cdot C/\omega^2$ for $s$ from $E_j$. Normal jets of order $b$ along a
curve $C$ appearing in $B_j$ with multiplicity $b$ are sections of $(L-bC)|_C$.

\begin{proposition}\label{prop:curves}
Let $C$ be an irreducible horizontal curve. (a) If $C\not\subseteq B_j$, then
$(L-\mu_jF)\cdot C\ge 0$. (b) If $C$ occurs in $B_j$ with multiplicity $b$, then
$(L-bC-\mu_jF)\cdot C\ge 0$.
\end{proposition}

\begin{proof}
(a) We have $L-\mu_jF=N_j+B_j+V_j$ with $N_j\cdot C\ge 0$, $B_j\cdot C\ge 0$, and
$V_j\cdot C\ge 0$. (b) Similarly, $L-bC-\mu_jF=N_j+(B_j-bC)+V_j$, and $C$ is not a
component of $B_j-bC$.
\end{proof}

\subsection{Source Bundles Not Pulled Back from the Base}
The identity of Theorem~\ref{thm:rankr} admits a generalization to an arbitrary
source bundle.

\begin{proposition}\label{prop:universal}
Let $\cF$ be locally free of rank $r\ge 2$ on $X$ and $\cF\to L$ a morphism with
image $A\otimes\cI_Z$, where $A$ is a line bundle and $Z$ has finite length. Then
the kernel $K$ is locally free of rank $r-1$, and
\[
\Bigl(A-\frac{c_1(\cF)}{r}\Bigr)^2=\frac{\Delta(K)+2(r-1)\ell(Z)}{r}-\frac{r-1}{r^2}\,\Delta(\cF).
\]
\end{proposition}

\begin{proof}
We have $c_1(K)=c_1(\cF)-A$ and $\ch_2(K)=\ch_2(\cF)-\frac12A^2+\ell(Z)$.
Inserting these into $\Delta(K)=c_1(K)^2-2(r-1)\ch_2(K)$, using
$\ch_2(\cF)=(c_1(\cF)^2-\Delta(\cF))/(2r)$, and collecting terms gives the
identity.
\end{proof}

For $\cF=\phi^*E$ one has $\Delta(\cF)=0$ and $c_1(\cF)\equiv(\deg E)F$, and the
proposition reduces to Theorem~\ref{thm:rankr}. When $\Delta(K)\ge 0$, the only
term that can make the left-hand side negative is $\Delta(\cF)>0$, which requires
a bundle that is not pulled back from the base. In the arithmetic setting, this
points back to Parshin's program of an arithmetic Bogomolov--Miyaoka--Yau
inequality~\cite{Pa88,MB90} and to arithmetic inequalities of Bogomolov--Gieseker
type~\cite{Mo94}. Soul\'e~\cite{So05} followed Miyaoka's approach through rank-two
extensions; he notes that the arithmetic second Chern number ``contains an
archimedean summand -- an integral over the set of complex points of $X$ -- which
is not easy to bound from above'', and he obtains lower bounds for some successive
minima rather than an upper bound for $\ob^{\,2}$. We draw no further conclusion
from this.

\section{Concluding Remarks}\label{sec:conclusion}

Within the class of arguments considered here, all available positivity is of one
kind: the intersection product of a nef divisor with an effective divisor is
nonnegative. In every case treated above, this positivity enters the leading
coefficient with the sign that makes the margin nonpositive, whereas a PMB
argument requires a positive margin. For pencils, this is established by
Theorem~\ref{thm:rank2} in rank two and by Corollary~\ref{cor:ss} for semistable
$\phi_*L$ in any rank. What remains open in this direction is the arithmetic
anomaly question for unstable $\phi_*L$ of rank $r\ge 3$ (Open
problem~\ref{prob:anomaly}). For~\cite{So26}, version~1, the consequence is
limited and concrete: Proposition~9\,vi) is false, inequality~(36) has a
negative leading coefficient, and the proof of Theorem~1 given there is
incomplete.

\phantomsection
\addcontentsline{toc}{section}{AI Disclosure}
\begingroup
\sloppy\emergencystretch=2em
% =============================================================================
% AI DISCLOSURE -- PIPELINE STANDARD, ENGLISH VERSION (since 2026-08-17, T-20260817-02)
% =============================================================================
% Language-pure version for English-only manuscripts. Text copied verbatim
% from the English block of ai_disclosure_STANDARD.tex (no content change) --
% only the heading is now language-pure and the German block is dropped. The
% bilingual file ai_disclosure_STANDARD.tex remains unchanged (existing
% references may keep pointing to it).
%
% DEVIATING DOWNWARD: If a paper genuinely used substantially less AI, do NOT
% switch to a different file -- weaken the wording here instead, e.g. replace
% "extensively in all phases" with the phases that actually applied, and
% strike whatever did not happen from the quality-assurance list. Principle:
% state honestly what happened -- do not hedge.
%
% Usage (English-only manuscripts only):
%   \input{../../_templates/disclosure/ai_disclosure_STANDARD_en}
% =============================================================================

\section*{AI Disclosure}

\noindent
AI systems were used extensively in all phases of this work: conception
and research discussion, literature and source work, mathematical analysis
and proof development, creation and checking of the computational scripts,
execution and evaluation of the computational runs, text, translation, and
typesetting.

\medskip

\noindent
Quality assurance: use of several independent models from different
providers with counter-checks and adversarial AI reviews; traceable,
versioned computational scripts with stored result files and checksums;
continuously maintained proof and register documents; strict claim hygiene
(assertions only with documented evidence, addenda instead of silent
corrections); changelogs across all versions.

\endgroup

\phantomsection
\addcontentsline{toc}{section}{References}
\begin{thebibliography}{99}
\small
\frenchspacing

\bibitem{BHPV04}
W.~P. Barth, K.~Hulek, C.~A.~M. Peters, A.~Van de Ven,
\emph{Compact Complex Surfaces}, 2nd enlarged ed.,
Ergebnisse der Mathematik und ihrer Grenzgebiete (3) 4, Springer, Berlin, 2004.

\bibitem{Bo78}
F.~A. Bogomolov,
Holomorphic tensors and vector bundles on projective varieties,
\emph{Izv. Akad. Nauk SSSR Ser. Mat.} 42 (1978), 1227--1287;
English transl.\ \emph{Math. USSR-Izv.} 13 (1979), 499--555.

\bibitem{CH88}
M.~Cornalba, J.~Harris,
Divisor classes associated to families of stable varieties, with applications
to the moduli space of curves,
\emph{Ann. Sci. \'Ec. Norm. Sup\'er.} (4) 21 (1988), 455--475.

\bibitem{HN75}
G.~Harder, M.~S. Narasimhan,
On the cohomology groups of moduli spaces of vector bundles on curves,
\emph{Math. Ann.} 212 (1975), 215--248.

\bibitem{HL10}
D.~Huybrechts, M.~Lehn,
\emph{The Geometry of Moduli Spaces of Sheaves}, 2nd ed.,
Cambridge Mathematical Library, Cambridge University Press, Cambridge, 2010.

\bibitem{Ko67}
K.~Kodaira,
A certain type of irregular algebraic surfaces,
\emph{J. Analyse Math.} 19 (1967), 207--215.

\bibitem{MD85}
M.~Martin-Deschamps,
La construction de Kodaira--Parshin,
in: \emph{S\'eminaire sur les pinceaux arithm\'etiques: la conjecture de Mordell},
Ast\'erisque 127 (1985), 261--273.

\bibitem{Mi77}
Y.~Miyaoka,
On the Chern numbers of surfaces of general type,
\emph{Invent. Math.} 42 (1977), 225--237.

\bibitem{MB90}
L.~Moret-Bailly,
Hauteurs et classes de Chern sur les surfaces arithm\'etiques,
in: \emph{S\'eminaire sur les pinceaux de courbes elliptiques (\`a la recherche
de \guillemotleft Mordell effectif\guillemotright)}, Ast\'erisque 183 (1990), 37--58.

\bibitem{Mo94}
A.~Moriwaki,
Inequality of Bogomolov--Gieseker type on arithmetic surfaces,
\emph{Duke Math. J.} 74 (1994), 713--761.

\bibitem{Pa88}
A.~N. Parshin,
The Bogomolov--Miyaoka--Yau inequality for the arithmetical surfaces and its
applications,
in: \emph{S\'eminaire de Th\'eorie des Nombres, Paris 1986--87},
Progr. Math. 75, Birkh\"auser, Boston, 1988, 299--312.

\bibitem{So05}
C.~Soul\'e,
Semi-stable extensions on arithmetic surfaces,
arXiv:math/0505079 (2005); in: \emph{Moduli Spaces and Arithmetic Geometry
(Kyoto, 2004)}, Adv. Stud. Pure Math. 45 (2006), 283--295.

\bibitem{So26}
C.~Soul\'e,
On the $abc$ conjecture,
arXiv:2611.05217, version~1 (3 November 2026).

\bibitem{Xi87}
G.~Xiao,
Fibered algebraic surfaces with low slope,
\emph{Math. Ann.} 276 (1987), 449--466.

\end{thebibliography}

\end{document}
